\section{Noncentered contracts and matched-accuracy implementation}\label{sec:r12-experiment}
We test the implementation consequences of the balance and acceptance-face results, not only whether a floating-point optimizer terminates. The current experiment complements, rather than replaces, the exact canonical economic comparison and the earlier nonlinear transfer study. Every adverse predecessor outcome remains in EC Section~\ref{sec:r10-nonlinear}.

\subsection{Design and reoptimized comparators}
The smooth resource objective is quartic maintenance plus quadratic temporal and spatial coupling,
\[
 \mathcal R_0(x)=\sum_n w_n\left\{m\|x_n\|^2+\gamma\sum_jx_{nj}^4
               +\lambda_s\|x_n-x_{\mathrm{pa}(n)}\|^2
               +\zeta\sum_{(i,j)\in\mathcal E}\omega_{ij}(x_{ni}-x_{nj})^2\right\}.
\]
Here $m=3/2$, $\lambda_s=3/5$, discount is $9/10$, and the initial installed tier is one half. The subscript on $\lambda_s$ distinguishes smooth adjustment from the absolute-switching parameter of Section~\ref{sec:r12-network}. The current computational study sets that absolute-switching term to zero; the full nonsmooth theorem is instead checked by independently formulated linear programs and exact examples. Net customer tariffs are $a_{nj}=1+(j\bmod3)/4$. Provider reward uses forcing $p=a+r$, where $r$ is an externally specified linear resource saving. A fixed resource offset of 100 per service--review can make resource accounting nonnegative on the generated box; it cancels from all reported gains and derivatives. This experiment does not endogenize prices or change the demand occupation law.

Four fixed graph families contain twelve independently generated instances each: three-service, three-review cycles; four-service, four-review random-edge graphs; six-service, four-review geometric graphs; and eight-service, four-review weighted-block graphs. Thus contingent plans have 21, 60, 90, and 120 coordinates. Parameters, graph generators, and random seeds are recorded before candidate evaluation. Replication identifiers are reused across families, forming common-random-number blocks; these are resampled jointly, not treated as 48 independent draws from one population. The latter two families have quartic curvature two and three and spatial coupling one half and one, respectively. Forcing is generated from heterogeneous rational service levels plus uncentered node shocks; it is not normalized to make one half optimal.

For each instance we optimize a separate constant-in-time tier for every service, including initial installation and spatial coupling. This vector static class is stronger than the single common-tier benchmark used in the original centered study. We round the optimizer to rational tiers and independently bound its static-class loss. The largest unnormalized static gap is \RtwStaticGap. Every subtree cap is then defined from that recorded outside protocol, without reoptimizing the cap for the adaptive method. Fifteen of the 48 outside protocols have at least one boundary tier, and the recorded tiers range from zero to one. Reported adaptive gains are relative to the rational outside protocol; its tiny static approximation bound is disclosed rather than described as exact equality with an irrational optimizer.

The first four previously frozen R10 training seeds supply both scalar-compatible critics, with no retraining on the new instances. They use the same 63 softplus features and the same scalar observations as before. The known static value and actor are re-anchored to the new comparator, but the fitted readout coefficients do not change. This is a deliberate transfer test. It is neither a claim that these four seeds exhaust training variability nor a test of an oracle that was trained on the shifted problems. The old ``value-only'' fit means no derivative targets in its loss; its shared feature scaling uses derivative geometry and is not derivative-free preprocessing.

\subsection{Repair, gate, and actual refinement}
We record raw box-clipped actors, all pre-repair continuation violations, the global scaling factor, and the exact reward effect of repair. A weighted Euclidean projection solves a different problem from radial scaling: it finds the closest feasible actor in the discounted tier metric, not the highest-reward actor. A separate tree-flow actor uses nonnegative child-to-parent payment transfers. Each transfer increases a parent tier by its flow divided by the parent's discounted tariff and decreases the child's tier by the flow divided by the child's discounted tariff. Every nonroot subtree payment therefore changes by the negative incoming flow, and root payment is unchanged. A common box-limited step preserves feasibility. These three approaches expose different parts of the acceptance geometry; none is assumed to identify the optimal exposed face for free.

The deployed learned pipeline generates an actor, projects, rounds and exactly repairs it, audits its reward and residual, and gates against the static protocol. It then actually runs SLSQP refinement when necessary. The cold classical pipeline starts from the same static protocol. Both use the identical quartic objective, analytic gradients, box bounds, every subtree constraint, and the same rational certificate with explicit upper- and lower-bound prices. The common audit schedule checks the initial candidate, every fifth iteration, and termination. Each pipeline stops only when the exact rational upper gap, divided by discounted service--reviews, is at most the prescribed tolerance. A numerical success flag cannot satisfy this stopping rule.

Table~\ref{tab:r12-matched} gives two tolerances and charges the learned pipeline for candidate generation, projection, audit, gating, and measured refinement. Comparator construction and common setup are recorded separately. All 432 new deployments meet the tighter tolerance, and their exactly evaluated final reward is at least the static reward. The maximum normalized final gap is \RtwFinalGap. An independent verifier imports neither the optimizer nor numerical linear-algebra libraries; across current and predecessor re-audits it checks \RtwCertificates\ candidate certificates, \RtwSubtrees\ subtree inequalities, and \RtwTiers\ tier inequalities. Five timing repetitions after a warmup, with randomized pipeline order, are reported separately in EC Table~\ref{tab:r12-repeated}; they are not treated as additional economic observations.

These matched-accuracy runs do not show a learned speed advantage over cold SLSQP. They establish an executable fallback that meets a stated accepted-value tolerance, including when transfer fails. A learned proposal's low generation time is not used to conceal the cost of obtaining the final certificate. We also execute the cached classical fallback on all forty original 992-coordinate weighted-block learned proposals. Its largest normalized final gap is \RtwLargeGap; the median measured generation--audit--fallback time is \RtwLargeMedianMs\ milliseconds. This is an explicit classical fallback, not a claim that those failed actors provide useful Newton warm starts.

\subsection{What the transfer evidence identifies}
EC Table~\ref{tab:r12-repair} reports every raw and repaired outcome. Of the 384 new learned proposals, \RtwRawNegative\ remain worse than static after weighted projection. Such failures are not omitted or relabeled as improvements. Projection can avoid some unnecessary scaling loss, but cannot generally correct a severely inaccurate critic. The tree-flow parameterization enforces the direction of continuation-compatible payment transfers; it does not replace optimization of their economic magnitudes. This agrees with Theorem~\ref{thm:r12-face}: feasibility and closeness alone do not establish the optimal-face condition or eliminate all first-order price losses.

The mean paired projection-gain difference between derivative-weighted and value-only actors is \RtwPairedMean\ per service--review. A descriptive crossed bootstrap, resampling training seeds and common replication blocks jointly across the four fixed graph families, has percentile interval $[\RtwPairedLower,\RtwPairedUpper]$. Its width and inclusion of zero preclude a reproducible superiority conclusion. This interval addresses different variation from R10's training-seed interval conditional on eight fixed problems. Neither is an inference over an unspecified population of graph families. The finite-sample selection theorem requires separately frozen pipelines, an independent validation distribution, and a valid uniform bound; it is not retroactively asserted by this stress test.

On the original weighted-block family, learned residual gaps are over a thousand times the classical gain for typical proposals. Using the static upper bound tightens the gate's error estimate but still leaves a large gap. EC Table~\ref{tab:r12-ratios} reports both ratios. Thus a no-harm gate is a feasibility/economic safeguard, not a substitute for accuracy refinement. The new endpoint-oracle study also varies the finite-difference step and certified value error; EC Table~\ref{tab:r12-fd} shows why derivative targets require an explicit error budget. The synthetic evidence supports audited implementation and diagnoses failure mechanisms, rather than supplying a field calibration or a universal architecture ranking.
\begin{table}[p]\centering
\caption{Certified refinement under noncentered forcing}\label{tab:r12-matched}
\setlength{\tabcolsep}{3.5pt}\begin{tabular}{llrrrrr}
\toprule
Family & Pipeline & \shortstack{Met\\tolerance} & \shortstack{Maximum\\final gap} & \shortstack{Mean\\final gain} & \shortstack{Time at\\$10^{-5}$} & \shortstack{Time at\\$10^{-7}$} \\
\midrule
Cycle & Cold SLSQP & 12/12 & 1.4e-08 & 0.004605 & 11.95 & 13.69 \\
Cycle & Value & 48/48 & 6.2e-08 & 0.004605 & 19.32 & 19.60 \\
Cycle & Value--gradient & 48/48 & 7.8e-08 & 0.004605 & 19.15 & 19.41 \\
Random-edge & Cold SLSQP & 12/12 & 8.3e-08 & 0.005110 & 30.90 & 35.78 \\
Random-edge & Value & 48/48 & 7.5e-08 & 0.005110 & 42.93 & 49.53 \\
Random-edge & Value--gradient & 48/48 & 8.8e-08 & 0.005110 & 43.30 & 49.91 \\
Geometric & Cold SLSQP & 12/12 & 9.4e-08 & 0.003731 & 54.24 & 67.84 \\
Geometric & Value & 48/48 & 9.7e-08 & 0.003731 & 73.06 & 88.62 \\
Geometric & Value--gradient & 48/48 & 9.7e-08 & 0.003731 & 73.54 & 86.91 \\
Weighted block & Cold SLSQP & 12/12 & 8.8e-08 & 0.001923 & 83.48 & 113.03 \\
Weighted block & Value & 48/48 & 8.8e-08 & 0.001923 & 110.87 & 140.69 \\
Weighted block & Value--gradient & 48/48 & 8.8e-08 & 0.001923 & 110.68 & 139.06 \\
\bottomrule
\end{tabular}
\begin{minipage}{.98\textwidth}\small
Gaps and gains are per discounted service--review. Times are median milliseconds across the declared instances and, for learned pipelines, frozen training seeds. They include candidate generation, projection, rational audit, gating, and actual refinement; a cold pipeline starts at the same optimized static protocol. Common setup is recorded separately. These are not independent timing repetitions or a speedup claim. The last column compares the same certified error tolerance, not different solver stopping flags.
\end{minipage}
\end{table}

